\documentclass{beamer}

\usepackage{amsmath,amsfonts,amssymb,bm}
\usepackage{quantikz}
\usepackage{graphicx}

\usetheme{Madrid}

\title{The HHL Algorithm}
\subtitle{Quantum Linear Systems and Spectral Methods}
\author{}
\date{}

\begin{document}

\frame{\titlepage}

%================================================
\section{Motivation}
%================================================

\begin{frame}{Linear Systems Everywhere}
\begin{itemize}
\item Solve:
\[
A \mathbf{x} = \mathbf{b}
\]
\item Appears in:
\begin{itemize}
\item PDEs
\item Machine learning
\item Many-body physics
\end{itemize}
\item Classical cost: $\mathcal{O}(N^3)$ (dense case)
\end{itemize}
\end{frame}

%------------------------------------------------

\begin{frame}{Quantum Goal}
\begin{itemize}
\item Encode vector as quantum state:
\[
|b\rangle = \sum_i b_i |i\rangle
\]
\item Goal:
\[
|x\rangle \propto A^{-1}|b\rangle
\]
\item Output is a quantum state, not a classical vector
\end{itemize}
\end{frame}

%================================================
\section{Spectral Decomposition}
%================================================

\begin{frame}{Eigenbasis Representation}
\[
A |u_j\rangle = \lambda_j |u_j\rangle
\]
\[
|b\rangle = \sum_j \beta_j |u_j\rangle
\]
\end{frame}

%------------------------------------------------

\begin{frame}{Formal Solution}
\[
A^{-1}|b\rangle = \sum_j \frac{\beta_j}{\lambda_j} |u_j\rangle
\]
\begin{itemize}
\item Problem reduces to eigenvalue inversion
\end{itemize}
\end{frame}

%================================================
\section{Quantum Phase Estimation}
%================================================

\begin{frame}{Hamiltonian Encoding}
\begin{itemize}
\item Assume $A$ is Hermitian
\item Implement:
\[
U = e^{iAt}
\]
\end{itemize}
\end{frame}

%------------------------------------------------

\begin{frame}{QPE Action}
\[
|u_j\rangle \rightarrow |u_j\rangle |\lambda_j\rangle
\]
\end{frame}

%------------------------------------------------

\begin{frame}{State After QPE}
\[
\sum_j \beta_j |u_j\rangle |\lambda_j\rangle
\]
\end{frame}

%================================================
\section{Controlled Rotation}
%================================================

\begin{frame}{Key Step: Inversion}
\begin{itemize}
\item Introduce ancilla qubit
\item Perform controlled rotation:
\[
|\lambda_j\rangle \rightarrow
\sqrt{1-\frac{C^2}{\lambda_j^2}}|0\rangle +
\frac{C}{\lambda_j}|1\rangle
\]
\end{itemize}
\end{frame}

%------------------------------------------------

\begin{frame}{State After Rotation}
\[
\sum_j \beta_j |u_j\rangle |\lambda_j\rangle
\left(
\sqrt{1-\frac{C^2}{\lambda_j^2}}|0\rangle +
\frac{C}{\lambda_j}|1\rangle
\right)
\]
\end{frame}

%================================================
\section{Uncomputation}
%================================================

\begin{frame}{Inverse QPE}
\begin{itemize}
\item Remove eigenvalue register
\item Leaves:
\[
\sum_j \frac{\beta_j}{\lambda_j}|u_j\rangle
\]
\end{itemize}
\end{frame}

%------------------------------------------------

\begin{frame}{Final State}
\[
|x\rangle \propto A^{-1}|b\rangle
\]
\end{frame}

%================================================
\section{Full Circuit}
%================================================

\begin{frame}{HHL Circuit (Overview)}
\centering
\begin{quantikz}
\lstick{$|b\rangle$} & \qw & \qw & \qw & \qw & \qw \\
\lstick{$|0\rangle^{\otimes m}$} & \gate{QPE} & \qw & \gate{QPE^\dagger} & \qw & \qw \\
\lstick{$|0\rangle$} & \qw & \gate{R(\lambda^{-1})} & \qw & \meter{} & \qw
\end{quantikz}
\end{frame}

%------------------------------------------------

\begin{frame}{Detailed Circuit Step 1: QPE}
\centering
\begin{quantikz}
\lstick{$|b\rangle$} & \ctrl{1} & \qw \\
\lstick{$|0\rangle$} & \gate{e^{iAt}} & \qw
\end{quantikz}
\end{frame}

%------------------------------------------------

\begin{frame}{Detailed Circuit Step 2: Controlled Rotation}
\centering
\begin{quantikz}
\lstick{$|\lambda\rangle$} & \ctrl{1} & \qw \\
\lstick{$|0\rangle$} & \gate{R_y(\lambda^{-1})} & \qw
\end{quantikz}
\end{frame}

%------------------------------------------------

\begin{frame}{Detailed Circuit Step 3: Uncomputation}
\centering
\begin{quantikz}
\lstick{} & \gate{QPE^\dagger} & \qw
\end{quantikz}
\end{frame}

%================================================
\section{Measurement}
%================================================

\begin{frame}{Post-selection}
\begin{itemize}
\item Measure ancilla
\item Condition on $|1\rangle$
\item Success probability depends on $\kappa$
\end{itemize}
\end{frame}

%================================================
\section{Complexity}
%================================================

\begin{frame}{Runtime}
\[
\mathcal{O}\left(\kappa^2 \log N\right)
\]
\begin{itemize}
\item $\kappa$ = condition number
\end{itemize}
\end{frame}

%------------------------------------------------

\begin{frame}{Conditions for Speedup}
\begin{itemize}
\item Sparse matrix
\item Efficient Hamiltonian simulation
\item Efficient state preparation
\end{itemize}
\end{frame}

%================================================
\section{Limitations}
%================================================

\begin{frame}{Readout Problem}
\begin{itemize}
\item Cannot extract full vector efficiently
\item Only expectation values accessible
\end{itemize}
\end{frame}

%------------------------------------------------

\begin{frame}{Condition Number Issue}
\begin{itemize}
\item Runtime scales with $\kappa$
\item Ill-conditioned systems are hard
\end{itemize}
\end{frame}

%================================================
\section{Physics Interpretation}
%================================================

\begin{frame}{Operator View}
\[
A^{-1} = \sum_j \frac{1}{\lambda_j} |u_j\rangle \langle u_j|
\]
\end{frame}

%------------------------------------------------

\begin{frame}{Connection to Green's Functions}
\[
G = (A)^{-1}
\]
\begin{itemize}
\item Resolvent operator
\item Linear response theory
\end{itemize}
\end{frame}

%------------------------------------------------

\begin{frame}{Many-Body Perspective}
\begin{itemize}
\item HHL = spectral filtering
\item Similar to:
\begin{itemize}
\item Imaginary time evolution
\item Projection methods
\end{itemize}
\end{itemize}
\end{frame}

%================================================
\section{Applications}
%================================================

\begin{frame}{Applications}
\begin{itemize}
\item Quantum machine learning
\item PDE solvers
\item Data fitting
\end{itemize}
\end{frame}

%================================================
\section{Summary}
%================================================

\begin{frame}{Summary}
\begin{itemize}
\item HHL solves $A x = b$ in quantum form
\item Uses:
\begin{itemize}
\item Phase estimation
\item Controlled rotation
\item Uncomputation
\end{itemize}
\item Connects strongly to spectral physics
\end{itemize}
\end{frame}

%================================================
\section{HHL and Green's Function Formalism}
%================================================

\begin{frame}{From Linear Systems to Resolvents}
A central observation is that HHL solves a linear system
\[
A x = b
\]
by preparing a quantum state proportional to
\[
|x\rangle \propto A^{-1}|b\rangle.
\]
In many-body physics, the inverse of an operator is precisely the basic structure of a
\emph{resolvent} or \emph{Green's function}.

If we define
\[
G(z) = (z I - H)^{-1},
\]
then solving
\[
(zI-H)|x\rangle = |b\rangle
\]
is equivalent to preparing
\[
|x\rangle = G(z)|b\rangle.
\]

Thus HHL may be viewed as a quantum algorithm for applying a resolvent operator to a source state.
\end{frame}

%------------------------------------------------

\begin{frame}{Spectral Form of the Green's Function}
Let
\[
H|u_j\rangle = E_j |u_j\rangle.
\]
Then the resolvent has the spectral decomposition
\[
G(z) = (zI-H)^{-1}
= \sum_j \frac{1}{z-E_j} |u_j\rangle \langle u_j|.
\]
If
\[
|b\rangle = \sum_j \beta_j |u_j\rangle,
\]
then
\[
G(z)|b\rangle = \sum_j \frac{\beta_j}{z-E_j}|u_j\rangle.
\]

This is structurally identical to the HHL transformation:
\[
A^{-1}|b\rangle
=
\sum_j \frac{\beta_j}{\lambda_j}|u_j\rangle.
\]

Hence, HHL implements the same spectral inversion mechanism that underlies Green's functions.
\end{frame}

%------------------------------------------------~
\begin{frame}{Physical Meaning of the Spectral Filter}
  The factor
  \[
  \frac{1}{z-E_j}
  \]
  acts as a spectral filter.

  \begin{itemize}
  \item Eigenmodes with energies near \(z\) are enhanced.
  \item Eigenmodes far from \(z\) are suppressed.
  \item The inverse operator encodes propagation and response properties.
  \end{itemize}

  In this sense, HHL is not merely a linear solver. It is a \emph{spectral transformation algorithm}:
  it maps a source state to a response state through an inverse operator.
\end{frame}

%------------------------------------------------

\begin{frame}{Retarded Green's Function}
  In condensed matter and many-body theory, one often uses
  \[
  G^R(\omega) = \frac{1}{\omega + i\eta - H},
  \qquad \eta > 0.
  \]
  The small positive \(\eta\) regularizes the inverse and enforces causal boundary conditions.

  Applying this operator to a source state gives
  \[
  |x(\omega)\rangle = G^R(\omega)|b\rangle.
  \]

  Formally, an HHL-type algorithm can be interpreted as a way of preparing a quantum state proportional to this frequency-dependent response vector, provided one can encode the shifted operator
  \[
  A(\omega) = \omega I - H
  \]
  or
  \[
  A(\omega) = \omega I + i\eta - H.
  \]
\end{frame}

%------------------------------------------------

\begin{frame}{Connection to Correlation Functions}
  Many observables in quantum many-body physics can be written in terms of Green's functions or resolvents.

  For an operator \(\hat{O}\), a typical linear-response quantity has the form
  \[
  \chi(\omega)
  =
  \langle \psi_0 |
  \hat{O}^\dagger
  \frac{1}{\omega + i\eta - (H-E_0)}
  \hat{O}
  |\psi_0\rangle.
  \]
  If we define the source state
  \[
  |b\rangle = \hat{O}|\psi_0\rangle,
  \]
  then
  \[
  \chi(\omega)
  =
  \langle b|
  \frac{1}{\omega + i\eta - (H-E_0)}
  |b\rangle.
  \]

  This shows that HHL-like inversion is directly connected to computing dynamical susceptibilities.
\end{frame}

%------------------------------------------------

\begin{frame}{HHL as a Green's Function Engine}
  At a conceptual level, HHL performs three tasks:
  \begin{enumerate}
  \item Decompose the source state into eigenmodes.
  \item Invert the spectral coefficients:
    \[
    \lambda_j \mapsto \frac{1}{\lambda_j}.
    \]
  \item Recombine the filtered eigenmodes into the response state.
  \end{enumerate}

  This is exactly what is needed for:
  \begin{itemize}
  \item resolvent methods,
  \item Green's function methods,
  \item frequency-domain linear response,
  \item projection-based many-body calculations.
  \end{itemize}
\end{frame}

%================================================
\section{HHL and Linear Response Theory}
%================================================

\begin{frame}{Linear Response in Physics}
  Suppose a system with Hamiltonian \(H_0\) is perturbed by
  \[
  V(t) = - f(t)\hat{F}.
  \]
  To first order, the expectation value of an observable \(\hat{O}\) changes as
  \[
  \delta \langle \hat{O}(t) \rangle
  =
  \int dt' \, \chi_{OF}(t-t') f(t'),
  \]
  where \(\chi_{OF}\) is the response function.

  In the frequency domain, one typically writes
  \[
  \delta \langle \hat{O}(\omega)\rangle
  =
  \chi_{OF}(\omega) f(\omega).
  \]

  The response function is often expressible through inverse operators of the form
  \[
  (\omega I - M)^{-1},
  \]
  which is precisely the structure targeted by HHL.
\end{frame}

%------------------------------------------------

\begin{frame}{Resolvent Form of the Susceptibility}
  A standard spectral representation is
  \[
  \chi_{OF}(\omega)
  =
  \sum_{n\neq 0}
  \left(
  \frac{\langle 0|\hat{O}|n\rangle \langle n|\hat{F}|0\rangle}
       {\omega-(E_n-E_0)+i\eta}
       -
       \frac{\langle 0|\hat{F}|n\rangle \langle n|\hat{O}|0\rangle}
            {\omega+(E_n-E_0)+i\eta}
            \right).
            \]
            This is a sum over poles, or equivalently a matrix element of a resolvent.

            Thus linear response is fundamentally an \emph{inverse-operator problem}, which places it in the same mathematical class as HHL.
\end{frame}

%================================================
\section{HHL and Time-Dependent Hartree--Fock}
%================================================

\begin{frame}{What is Time-Dependent Hartree--Fock?}
  Time-dependent Hartree--Fock (TDHF) describes the evolution of a Slater determinant under a time-dependent mean field.

  At the one-body density-matrix level,
  \[
  i\hbar \frac{d\rho}{dt} = [h[\rho],\rho].
  \]

  To study small oscillations around a stationary solution \(\rho_0\), one linearizes:
  \[
  \rho(t) = \rho_0 + \delta\rho(t),
  \]
  leading to an equation of motion for the fluctuation \(\delta\rho\).
\end{frame}

%------------------------------------------------

\begin{frame}{Linearized TDHF}
  After linearization, one obtains a matrix equation of the form
  \[
  i\hbar \frac{d}{dt}\delta \rho = \mathcal{L}\,\delta \rho + s(t),
  \]
  where:
  \begin{itemize}
  \item \(\mathcal{L}\) is the linearized Liouvillian or RPA matrix,
  \item \(s(t)\) is an external driving term.
  \end{itemize}

  In the frequency domain, this becomes
  \[
  (\omega I - \mathcal{L}) \delta \rho(\omega) = s(\omega).
  \]

  This is now explicitly a linear system. Its solution is
  \[
  \delta \rho(\omega) = (\omega I - \mathcal{L})^{-1} s(\omega).
  \]

  That is exactly the algebraic structure addressed by HHL.
\end{frame}

%------------------------------------------------

\begin{frame}{TDHF as an HHL-Type Problem}
  The formal similarity is immediate:
  \[
  A x = b
  \qquad \Longleftrightarrow \qquad
  (\omega I - \mathcal{L}) \delta \rho(\omega) = s(\omega).
  \]
  Hence:
  \[
  A \leftrightarrow \omega I - \mathcal{L},
  \qquad
  x \leftrightarrow \delta \rho(\omega),
  \qquad
  b \leftrightarrow s(\omega).
  \]

  If one can encode the linear-response matrix \(\mathcal{L}\) into a quantum representation, then an HHL-type algorithm could in principle prepare a state proportional to the TDHF response.
\end{frame}

%------------------------------------------------

\begin{frame}{Connection to Random Phase Approximation (RPA)}
  The small-amplitude limit of TDHF is closely related to the Random Phase Approximation (RPA).

  One typically obtains a non-Hermitian eigenvalue problem of the form
  \[
  \begin{pmatrix}
    A & B \\
    -B^* & -A^*
  \end{pmatrix}
  \begin{pmatrix}
    X \\ Y
  \end{pmatrix}
  =
  \omega
  \begin{pmatrix}
    X \\ Y
  \end{pmatrix}.
  \]
  Equivalently, one can formulate a driven linear-response problem
  \[
  (\omega I - M)\,v = s.
  \]

  Again, this is an inverse problem. The response amplitudes are obtained by applying the resolvent
  \[
  (\omega I - M)^{-1}
  \]
  to a source vector.
\end{frame}

%------------------------------------------------

\begin{frame}{What HHL Would Compute in TDHF}
  In a TDHF or RPA context, one is often interested in:
  \begin{itemize}
  \item transition amplitudes,
  \item density oscillations,
  \item response to an external field,
  \item strength functions.
  \end{itemize}

  These quantities can often be expressed as matrix elements involving the response vector:
  \[
  S(\omega) \sim \langle s | (\omega I - M)^{-1} | s \rangle.
  \]

  Thus an HHL-type routine would not need to output the full vector \(\delta\rho(\omega)\). It would already be useful if it allowed efficient estimation of such expectation values.
\end{frame}

%------------------------------------------------

\begin{frame}{Strength Functions and Dynamical Response}
  A common observable in TDHF and RPA is the strength function
  \[
  S_F(\omega)
  =
  -\frac{1}{\pi}\operatorname{Im}\chi_F(\omega),
  \]
  where
  \[
  \chi_F(\omega)
  =
  \langle 0|\hat{F}^\dagger
  \frac{1}{\omega + i\eta - (H-E_0)}
  \hat{F}|0\rangle.
  \]

  This is exactly a Green's function expression. Hence:
  \[
  \text{TDHF/RPA response}
  \quad \longleftrightarrow \quad
  \text{resolvent evaluation}
  \quad \longleftrightarrow \quad
  \text{HHL-type inversion}.
  \]
\end{frame}

%------------------------------------------------

\begin{frame}{Why This Connection is Interesting}
  The HHL framework is potentially relevant because:
  \begin{itemize}
  \item linear response is naturally formulated as a linear system,
  \item inverse operators define propagators and susceptibilities,
  \item many observables are expectation values rather than full vectors,
  \item this aligns well with quantum output models.
  \end{itemize}

  From a physics viewpoint, HHL is best interpreted not only as a solver of \(Ax=b\), but as a quantum method for preparing \emph{response states} and probing \emph{spectral properties}.
\end{frame}

%------------------------------------------------

\begin{frame}{Important Caveats}
  There are, however, important obstacles:
  \begin{itemize}
  \item The matrix \(A\) must be efficiently block-encoded or simulated.
  \item TDHF/RPA matrices may be non-Hermitian.
  \item Condition numbers may be large near resonances.
  \item State preparation and measurement may dominate the cost.
  \end{itemize}

  Therefore, the connection is mathematically deep and physically appealing, but practical advantage depends strongly on encoding, sparsity, conditioning, and the observables of interest.
\end{frame}

%------------------------------------------------

\begin{frame}{Hermitian Embedding of Non-Hermitian Problems}
  If the response matrix \(M\) is non-Hermitian, one common strategy is Hermitian embedding:
  \[
  \mathcal{H}
  =
  \begin{pmatrix}
    0 & M \\
    M^\dagger & 0
  \end{pmatrix}.
  \]
  Then \(\mathcal{H}\) is Hermitian, and one may reformulate the problem in an enlarged space.

  This is often the natural bridge between:
  \begin{itemize}
  \item non-Hermitian response equations,
  \item quantum phase estimation or HHL-type methods,
  \item Green's function formulations on quantum hardware.
  \end{itemize}
\end{frame}

%------------------------------------------------

\begin{frame}{Conceptual Summary}
  \begin{block}{Key idea}
    HHL applies an inverse operator spectrally:
    \[
    A^{-1} = \sum_j \frac{1}{\lambda_j}|u_j\rangle\langle u_j|.
    \]
  \end{block}

  \begin{block}{Green's function viewpoint}
    Choose
    \[
    A = zI - H,
    \]
    then HHL prepares a state proportional to
    \[
    (zI-H)^{-1}|b\rangle.
    \]
  \end{block}

  \begin{block}{TDHF / linear-response viewpoint}
    Choose
    \[
    A = \omega I - \mathcal{L},
    \]
    then HHL prepares a state proportional to the driven response
    \[
    \delta\rho(\omega) = (\omega I - \mathcal{L})^{-1} s(\omega).
    \]
  \end{block}
  \end{frame}


\end{document}



